\section{Introduction}

Large Language Models have achieved remarkable success through the AR paradigm, modeling sequences via next-token prediction with strict left-to-right causal dependencies~\citep{hurst2024gpt4o,grattafiori2024llama3,yang2025qwen3}. This approach naturally aligns with the sequential structure of language and enables efficient training through next-token likelihood maximization. However, the very success of this paradigm creates fundamental limitations: the sequential generation process imposes severe inference bottlenecks, precluding parallelization, and increasing latency at scale, while the rigid causal structure can be suboptimal for tasks requiring bidirectional reasoning and holistic understanding.

Discrete Masked Diffusion Language Models (MDLM) have emerged as a compelling alternative to the prevailing AR paradigm. By reconstructing sequences from random masked inputs, these models inherently support parallel generation and leverage a full bidirectional context, offering a different architectural approach~\citep{diffullama,yu_discrete_2025}. Although these conceptual advantages are clear, the field is still in an early developmental stage. Current research is actively focused on key challenges, including the refinement of specialized training regimes, the design of efficient sampling strategies, the efficient inference of opensource models, and reinforcement learning for MDLM. As a result of this ongoing exploration, most existing diffusion models, including recent advancements like Block Diffusion Language Models (BDLMs)~\citep{Blockdiffusion2025}, operate at a smaller scale (\textit{e.g.}, $\leq$8B parameters). Bridging this scale difference to the hundreds of billions of parameters seen in the leading mainstream AR models is a primary frontier for enabling diffusion models to fully capture complex linguistic patterns for practical deployment.

In this work, we introduce LLaDA2.0 series with 100B/16B total parameters diffusion language models that resolves these fundamental challenges through a novel two-stage continual pre-training (CPT) paradigm. Rather than attempting to train diffusion models from scratch, we leverage existing AR checkpoints as the foundation for a systematic conversion process that preserves linguistic knowledge while introducing diffusion capabilities.


The first stage, CPT aims to transform the foundational AR model into a capable diffusion language model. However, direct conversion is challenging due to the inherent data distribution gap between left-to-right generation and bidirectional denoising. Although the BDLM formulation partially reduces this gap through blockwise masked reconstruction, it suffers from low data utilization, limiting the effective exploitation of large-scale corpora. To this end, we introduce the Warmup–Stable–Decay (WSD) strategy, smoothly bridging the AR-to-dLLM gap while substantially improving CPT efficiency. WSD gradually expands the model’s receptive field to introduce diffusion-style context (\textit{Warmup}), strengthens global denoising under full-sequence training (\textit{Stable}), and then refines the model into an efficient blockwise structure (\textit{Decay}). This progressive adjustment enables a stable and data-efficient transition to diffusion-based learning. Additionally, under full attention in packed training sequences, diffusion models risk forming spurious dependencies across document boundaries, leading to semantic confusion and instability in bidirectional training. To prevent such cross-document interference, we introduce a document-level attention mask that restricts self-attention within individual documents, ensuring coherent context modeling.


The second stage, Post-training for Practical Deployment, transitions the model from a raw predictive engine into a capable and efficient assistant. The random masking nature of the diffusion fine-tuning objective means any single sample provides only a partial learning signal. We address this by employing a complementary masking strategy, which ensures near-100\% data utilization and accelerates convergence by guaranteeing every token contributes to the model's learning. With an efficient foundation for instruction tuning, we then align the model with human preferences by adapting modern techniques like Direct Preference Optimization (DPO)—originally designed for AR models—by reformulating the objective over the model's reconstruction loss. Beyond alignment, practical deployment hinges on inference speed. To realize the full promise of parallel decoding, which is often limited by a model's lack of predictive confidence, we incorporate an auxiliary confidence prediction loss. This trains the model to be ``sharper" and more certain, unlocking aggressive and efficient parallel generation without degrading quality.


We release instruction-tuned variants for practical deployment: \texttt{LLaDA2.0-mini} (16B parameters) for resource-constrained applications and \texttt{LLaDA2.0-flash} (100B parameters) for high-performance scenarios. Both variants retain the parallel decoding advantages of our diffusion training while being optimized for instruction following and safety through comprehensive post-training alignment.

Our contributions provide a practical recipe for the community to leverage AR stability while achieving diffusion parallelism, opening new possibilities for efficient large-scale language modeling.
